\section{总结与展望}

\subsection{课程设计工作总结}

本次课程设计围绕"视觉退化条件下 EEG 辅助视觉语言模型进行语义识别与图像描述生成"这一任务，在 Thought2Text 数据集上完成了从方案设计、代码实现、参数调优到效果评估的全流程实践。取得的阶段性成果如下：

\textbf{(1) 搭建了面向视觉退化场景的 EEG 辅助语义识别实验体系。}定义了六类退化条件（clean、lowres16、strong\_blur、strong\_noise、occlusion50、mixed）与五种 EEG 对照模式（real EEG、vision-only、shuffled EEG（打乱脑电）、random EEG、eeg-only），并采用 image-level 划分方式避免数据泄漏。

\textbf{(2) 实现了 A2 时间-频谱-空间 EEG 语义融合模型。}该模型在 CLIP 语义空间中以三路并行编码加门控残差融合的方式运行。在当前数据集与实验配置下，real EEG 在退化图像语义识别任务上的表现优于 vision-only 及 shuffled/random EEG 对照组。在全部 18 个 seed-condition 组合中，real EEG 均领先于所有对照设置。

\textbf{(3) 完成了 VTF token 级 EEG 视觉特征增强。}将 EEG 信息的注入位置从 pooled embedding 层面推进至 visual token 层面，为生成式视觉语言模型提供了 EEG 增强的视觉输入表示。

\textbf{(4) 构建了基于 Qwen2-VL 的生成式图像描述原型系统。}借助 LoRA 低秩适配与 best-of-N 重排序策略，在 strong 退化条件下取得了 97.8\% 的有效输出率，且 real EEG 相比 vision-only 在 caption class-hit 指标上提升了约 8.4 个百分点。

\subsection{主要收获}

通过本次课程设计，我们在 EEG 信号分析、多模态融合架构、预训练视觉语言模型的适配与微调、LoRA 参数高效训练方法以及严格的对照实验设计等方面均获得了较为系统的实操经验。整个开发过程经历了从最初的概念验证，到逐步扩展功能模块，再到从分类任务过渡至生成式输出的递进路线。每个阶段都产出可验证的中间结果，这种分层递进的构建方式对理解复杂系统的开发流程具有较好的参考价值。

\subsection{不足与改进方向}

本设计仍有若干方面值得进一步完善：(1) 数据集规模偏小——Thought2Text 仅涵盖 40 个类别约 12000 条 trial，制约了模型在更广类别范围上的泛化能力。(2) 生成式描述的训练以 BLIP 生成的伪 caption 为监督信号，其质量与人工标注之间仍存在差距，可能对生成模型的能力上限构成限制。(3) 未对 EEG 个体差异进行显式建模，全部训练在 6 名被试的混合数据上完成。(4) VTF token fusion 的交互层级较浅，若采用更复杂的 multi-layer bridge 设计，有望进一步提升增强后视觉 token 的质量。

后续可探索的方向包括：在更大规模的 EEG-image-text 数据集上进行预训练与迁移学习；引入 subject adaptation 模块以削弱个体差异带来的影响；采用更强的 Q-Former 或 Perceiver bridge 替换当前的双层 cross-attention 结构来执行 token fusion。

最后需要指出，本课程设计是在小样本 EEG-image 数据条件下开展的初步探索。实验结果表明，真实的配对 EEG 在视觉退化场景下能够为视觉语言模型提供一定程度的语义辅助信息。但受数据规模与模型复杂度的限制，当前系统生成的自由文本描述质量与大规模视觉语言模型在清洁图像上的输出仍有较大差距。本设计的工作定位是通过严格的对照实验，为 EEG 作为视觉语义增强信号的可行性提供初步的经验性参考。我们构建的是一个\textbf{小样本 EEG-enhanced generative VLM 原型}，在工程与算法层面仍有大量改进空间。